\chapter*{Declaración de Ética}
\addcontentsline{toc}{chapter}{Declaración de Ética}
\markboth{Declaración de Ética}{Declaración de Ética}

El presente documento académico y de investigación ha sido elaborado con estricto apego a los principios deontológicos del hacking ético, las directrices de pruebas de seguridad de OWASP (\textit{Open Web Application Security Project}) y los estándares internacionales de ciberseguridad defensiva. 

Las metodologías expuestas, en particular el modelado de amenazas \textbf{STRIDE} (desarrollado por Microsoft) y la cadena de intrusión cibernética \textbf{Cyber Kill Chain} (formulada por Lockheed Martin), se aplican bajo un enfoque puramente analítico, arquitectónico y documental sobre la infraestructura de tecnologías de la información de la Universidad Nacional del Altiplano (UNAP) de Puno.

\textbf{Compromiso de No Agresión e Inocuidad:}
\begin{itemize}
    \item Todo comando, simulación o vector de ataque presentado a lo largo de este trabajo ha sido ejecutado en laboratorios aislados de prueba, sobre entornos controlados de simulación o empleando información estrictamente pública y accesible mediante técnicas de Inteligencia de Fuentes Abiertas (\textit{Open-Source Intelligence} - OSINT).
    \item \textbf{No se han realizado ataques destructivos, explotación de vulnerabilidades sin autorización, interceptación de tráfico de red real ni alteración de bases de datos o sistemas en producción de la UNAP.}
    \item Toda información sensible, tales como direcciones IP internas, hashes criptográficos o parámetros internos de bases de datos institucionales, ha sido anonimizada y ofuscada para preservar la confidencialidad y la privacidad institucional.
\end{itemize}

El propósito fundamental de este estudio es proporcionar a la Oficina de Tecnologías de la Información (OTI) de la UNAP una evaluación estructurada que sirva como insumo para el diseño de controles preventivos, detección temprana de intrusiones y mejora continua de la resiliencia operativa.
